\usepackage[utf8]{inputenc}
\usepackage[T1]{fontenc}
\usepackage{textcomp}
\usepackage[english]{babel}
\usepackage{amsmath, amssymb, amsthm}
\usepackage{hyperref}
\usepackage{cleveref}
\usepackage{geometry}

% figure support
\usepackage{import}
\usepackage{xifthen}
\pdfminorversion=7
\usepackage{pdfpages}
\usepackage{transparent}
\newcommand{\incfig}[1]{%
  \def\svgwidth{\columnwidth}
  \import{./Figures/}{#1.pdf_tex}
}

\usepackage[most]{tcolorbox}

\pdfsuppresswarningpagegroup=1

\newcommand{\R}{\mathbb{R}}
\newcommand{\Q}{\mathbb{Q}}
\newcommand{\Z}{\mathbb{Z}}
\newcommand{\N}{\mathbb{N}}
\renewcommand{\H}{\mathbb{H}}
\newcommand{\F}{\mathbb{F}}
\newcommand{\W}{\mathbb{W}}

\usepackage{tikz}
\usepackage{xparse}

% Define a new environment for polar grid
\NewDocumentEnvironment{polargrid}{O{4}O{4}O{12}O{4.5}}{%
  % #1 = max radius (default: 4)
  % #2 = number of concentric circles (default: 4)
  % #3 = number of rays (default: 12)
  % #4 = axis length beyond max radius (default: 4.5)
  \begin{tikzpicture}[>=latex]
    % Draw the outer circle
    \draw[lightgray, thin] (0,0) circle (#1);
    % Draw concentric circles
    \foreach \r in {1,...,\numexpr#2-1\relax} {
      \draw[lightgray, thin] (0,0) circle (\r*#1/#2);
    }
    % Draw rays
    \foreach \i in {0,...,\numexpr#3-1\relax} {
      \pgfmathsetmacro{\angle}{(\i*360)/#3}
      \draw[lightgray, thin] (\angle:#1) -- (\angle:0);
    }
    % Draw main axes
    \foreach \angle in {0,90,180,270} {
      \draw[black, thick] (\angle:#1) -- (\angle:0);
    }
    % Angle labels
    \foreach \i/\label in {
      30/$\frac{\pi}{6}$,
      60/$\frac{\pi}{3}$,
      120/$\frac{2\pi}{3}$,
      150/$\frac{5\pi}{6}$,
      180/$\pi$,
      210/$\frac{7\pi}{6}$,
      240/$\frac{4\pi}{3}$,
      270/$\frac{3\pi}{2}$,
      300/$\frac{5\pi}{3}$,
      330/$\frac{11\pi}{6}$
    } {
      \draw (\i:#1 + 0.3) node {\label};
    }
    % Draw axis labels
    \draw[->, thick] (0,0) -- (0,#4) node[above] {$\mathbf{y}$};
    \draw[->, thick] (0,0) -- (#4,0) node[right] {$\mathbf{x}$};
  }{%
  \end{tikzpicture}
}

% Command for plotting points on polar grid
% This version properly handles negative radii by adding 180 degrees to the angle
\NewDocumentCommand{\polarpoint}{mmO{red}O{}}{%
  % #1 = radius
  % #2 = angle in radians (will be converted to degrees)
  % #3 = point color (default: red)
  % #4 = label (optional)

  % Handle negative radius by adding 180° to angle and using absolute radius
  \pgfmathsetmacro{\radius}{abs(#1)}
  \pgfmathsetmacro{\angle}{#2*180/pi}

  % If radius is negative, adjust the angle by adding 180°
  \ifdim #1 pt < 0 pt
    \pgfmathsetmacro{\angle}{\angle + 180}
  \fi

  % Normalize angle to be between 0° and 360°
  \pgfmathsetmacro{\angle}{mod(\angle, 360)}

  \draw[fill=#3] (\angle:\radius) circle (2pt) 
  \ifx\relax#4\relax\else
    node[above right, font=\scriptsize] {#4}
  \fi;
}

% For convenience, a command that formats polar coordinates nicely
\NewDocumentCommand{\polarcoord}{mm}{%
  % #1 = radius
  % #2 = angle (in radians)
  $(#1, \formatangle{#2})$%
}

% Helper command to format angles nicely
\NewDocumentCommand{\formatangle}{m}{%
  \ensuremath{#1}%
}

\usepackage[most]{tcolorbox}
\tcbuselibrary{listings, breakable, skins}

% Problem box
\newtcolorbox[auto counter, number within=section]{problem}[1][]{%
  colback=red!5!white,
  colframe=red!100!black,
  fonttitle=\bfseries,
  title=Problem~\thetcbcounter,
  label=#1,
  enhanced,
  breakable,
  sharp corners,
  boxrule=0.8pt,
  coltitle=black,
  before upper={\parindent15pt}, % optional formatting
}

% Solution box
\newtcolorbox[auto counter, number within=section]{solution}[1][]{%
  colback=green!5!white,
  colframe=green!100!black,
  fonttitle=\bfseries,
  title=Solution~\thetcbcounter,
  label=#1,
  enhanced,
  breakable,
  sharp corners,
  boxrule=0.8pt,
  coltitle=black,
  before upper={\parindent15pt}, % optional formatting
}
